\section{Conclusions}
The final chemical-process basis remains the INL/NGNP high-temperature nuclear-assisted SMR+CCS case: 130 MMSCFD H$_2$, 871$^\circ$C reforming, 176.8 MWth nuclear process heat and 17.3 MWe process demand~\cite{inltev953,inltev961}. At 85\% availability the declared lifecycle boundary gives approximately 0.917 MtCO$_2$e/y avoided, well above the CN4252 annual threshold.

The corrected hardware/economic calculation uses one coherent Nishihara GTHTR300C source architecture rather than mixing JAEA variants. The source split is 370 MWth hydrogen heat plus a 230 MWth power branch producing 88 MWe; the project's 176.8 MWth process draw is mapped explicitly to the residual reactor branch. Under that screening mapping, the zero-electricity-value case remains below S\$100/tCO$_2$e, while electricity-value cases provide additional sensitivities.

The resulting joint pass is a \textbf{conditional model result pending independent reviewer re-verification}. It is not a claim of demonstrated Singapore commercial feasibility or a guaranteed electricity offtake. The INL process design, JAEA hardware/economic source case and project-derived integration quantities are kept explicitly distinct.
